% method.tex

\chapter{基于TimeMixer的改进模型设计}
\echapter{Improved TimeMixer Model Design}

\section{改进模型总体框架}
\esection{Overall Framework}

本文以 TimeMixer 的多尺度分解混合结构为基础，围绕下采样、序列分解和变量交互三个关键环节进行增强。该设计保留了原模型轻量化和多尺度建模的优势，同时便于在统一框架下比较不同结构模块对预测性能的影响。

本文将改进方向概括为三个模块。第一个模块是深度可分离卷积下采样，记为 A，用于替代原始模型中的平均池化下采样。第二个模块是多核门控分解，记为 B，用于替代单一固定窗口的滑动平均分解。第三个模块是通道混合，记为 C，用于在通道独立建模之前引入一次变量间信息交互。三个模块分别对应原始模型中的多尺度构建、序列分解和变量处理环节。
这三个环节分别影响模型面对时间序列时的不同阶段。多尺度构建决定了原始输入在进入主干网络前如何被压缩和组织，若下采样方式过于固定，局部变化可能在早期被平滑；序列分解决定了模型如何区分较平滑的趋势和剩余波动，若窗口选择不合适，趋势估计可能过粗或过细；变量处理则决定了多变量序列中不同变量是否存在信息交换，若完全通道独立，模型效率较高，但变量间关系利用不足。本文选择在这三个位置分别加入轻量模块，是因为它们对应 TimeMixer 中较明确的结构环节，便于在不改变主干预测框架的情况下观察单个设计的影响。

本文改进模型的整体流程可以概括为：首先根据下采样参数构建多尺度输入序列；随后对各尺度序列分别进行归一化，并在启用通道混合时，于通道独立拆分和嵌入之前进行一次浅层变量交互；接着在各尺度上进行趋势项和季节项分解，并通过 TimeMixer 原有的跨尺度混合结构传播不同尺度的信息；最后由未来多尺度预测融合模块得到输出结果。本文三个改进模块均嵌入在上述流程中，而不是改变预测任务的输入输出形式，因此 M0 与各改进变体之间具有较好的可比性。
这种设计方式也决定了本文采用全组合消融实验，而不是只比较 M0 与 M4。完整模型 M4 同时包含三个模块，能够反映结构最完整时的结果，但它无法说明具体收益来自哪一个环节，也无法判断不同模块之间是否存在相互干扰。通过 M1、M2、M3 可以分别观察单模块的独立作用；通过 M12、M13、M23 可以进一步观察两两组合是否产生互补；再与 M4 对照，则可以分析三个模块同时加入时是否带来线性叠加收益。这样的消融设计使后续实验结论更容易对应到具体结构，而不是只停留在整体误差变化上。

\begin{figure}[htbp]
    \centering
    \includegraphics[width=0.98\textwidth]{thesis/fig_improved_timemixer_framework.png}
    \caption{改进TimeMixer模型整体框架图}
    \label{fig:improved_framework}
\end{figure}

表 \ref{tab:module_code_mapping} 概括了三个改进模块在模型流程中的作用。三个模块分别对应多尺度构建、趋势/季节分解和多变量输入处理三个环节，后续消融实验将围绕这些模块展开。

\begin{table}[htbp]
    \centering
    \caption{改进模块功能对应关系}
    \label{tab:module_code_mapping}
    \resizebox{\textwidth}{!}{%
    \begin{tabular}{cccl}
        \toprule
        模块 & 对应环节 & 改进方式 & 主要作用 \\
        \midrule
        深度可分离卷积下采样 & 多尺度构建 & 可学习卷积下采样 & 提升下采样过程的局部特征提取能力 \\
        多核门控分解 & 趋势/季节分解 & 多窗口趋势估计与门控融合 & 增强分解模块对不同周期尺度的适应性 \\
        通道混合 & 多变量输入处理 & 低秩残差映射 & 在通道独立建模前引入变量间信息交互 \\
        \bottomrule
    \end{tabular}}
\end{table}

\section{深度可分离卷积下采样模块}
\esection{Depthwise Separable Convolution Down-sampling}

原始 TimeMixer 在构建多尺度序列时默认使用平均池化。平均池化可以表示为对局部窗口内时间点进行等权平均：
\begin{equation}
    y_t = \frac{1}{k}\sum_{i=0}^{k-1} x_{t+i},
\end{equation}
其中 $k$ 为池化窗口大小。该操作没有可学习参数，计算稳定，但无法根据数据特征自动调整不同时间点的重要性。
平均池化的优势在于稳定和简单。它不会额外增加模型参数，也不会引入新的优化目标，因此在许多平滑序列上能够起到抑制噪声的作用。但这种操作对窗口内每个时间点一视同仁，无法区分局部峰值、突变点或短时间异常波动是否具有预测价值。对于长序列预测而言，如果重要局部信息在下采样阶段已经被平滑，后续跨尺度混合模块即使具有较好的表达能力，也难以恢复这些细节。因此，本文将下采样视为可以改进的前置环节，而不是只在主干网络内部增加参数。

为缓解这一问题，本文设计了深度可分离卷积下采样模块。该模块输入张量形状为 $(B,C,T)$，其中 $B$ 为批大小，$C$ 为变量数，$T$ 为时间长度。模块结构为：
\begin{equation}
    \mathrm{DWConv1d} \rightarrow \mathrm{LayerNorm} \rightarrow \mathrm{GELU} \rightarrow \mathrm{PWConv1d}.
\end{equation}
其中，逐通道卷积（Depthwise Conv1d）在每个变量内部沿时间维度提取局部特征，并通过步长完成下采样；LayerNorm 用于稳定中间特征分布；GELU 激活提供非线性表达能力；逐点卷积（Pointwise Conv1d）使用 $1 \times 1$ 卷积对通道信息进行线性变换。
本文没有直接采用普通一维卷积替代平均池化，主要是考虑到多变量时间序列的变量数可能较多。普通卷积会同时在时间维度和通道维度上进行混合，参数量随输入通道和输出通道同时增加，在 ECL、Traffic 等高维数据集上容易带来额外开销。深度可分离卷积将这一过程拆成两步：Depthwise Conv1d 只在每个变量内部沿时间方向提取局部模式，关注短期变化和局部形态；Pointwise Conv1d 再对通道表示做线性映射，用较低成本调整特征空间。这样既保留了可学习下采样的能力，又避免把下采样模块设计得过重。

与普通卷积相比，深度可分离卷积将时间局部特征提取和通道映射拆开，参数量和计算量更低。与平均池化相比，它的卷积核可通过训练更新，因此下采样过程不再是固定平滑，而是具有一定任务适应性。本文并不认为该模块一定在所有数据集上优于平均池化，因为对于本身较平滑的数据，平均池化的平滑作用也可能有利。本文实验正是通过 M1 变体单独验证该模块的实际作用。
\begin{figure}[htbp]
    \centering
    \includegraphics[width=0.95\textwidth]{thesis/fig_depthwise_separable_downsampling.png}
    \caption{深度可分离卷积下采样模块结构图}
    \label{fig:depthwise_separable_downsampling}
\end{figure}

\section{多核门控分解模块}
\esection{Multi-kernel Gated Decomposition}

原始 TimeMixer 使用滑动平均分解序列，将输入拆分为季节项和趋势项。固定窗口滑动平均可以得到较稳定的趋势估计，但窗口大小一旦确定，就难以同时兼顾短期波动和长期趋势。实际数据中，不同数据集的周期结构差异明显，即使在同一数据集中，不同时间段的变化尺度也可能不同。
单一窗口分解隐含了一个假设：序列中适合用于趋势估计的时间尺度是固定的。但在真实时序数据中，这一假设往往较强。例如，交通序列可能同时包含小时级拥堵变化和日周期出行规律；气象序列可能在较短时间内出现突变，也可能在较长时间范围内体现稳定趋势；太阳能序列既受昼夜周期影响，也会受到局部天气扰动。若窗口过小，趋势项可能混入较多局部波动；若窗口过大，短周期结构又可能被过度平滑。多窗口分解的目的就是让模型同时获得几种不同平滑程度的趋势候选，再根据输入样本选择合适的融合方式。

为增强分解模块的适应能力，本文提出多核门控分解模块。该模块并行使用多个移动平均窗口，窗口大小为 $(3,11,37)$。对于输入序列 $\mathbf{X}$，首先通过三个不同窗口得到趋势候选：
在本文设置中，$k=3$ 可以理解为较短窗口，主要保留临近时间点的局部变化；$k=11$ 对应中等尺度窗口，能够在一定程度上平滑短期噪声，同时保留部分周期形态；$k=37$ 则产生更平滑的趋势估计，更偏向捕捉较长时间范围内的变化方向。三个窗口并不是对具体数据周期的精确划分，而是为模型提供从局部到平滑的多种趋势候选。相比只调整单一窗口大小，这种并行结构能够避免模型被固定在某一个时间尺度上。
\begin{equation}
    \mathbf{T}_j = \mathrm{AvgPool}_{k_j}(\mathbf{X}), \quad k_j \in \{3,11,37\}.
\end{equation}
此处 $\mathrm{AvgPool}$ 表示步长为 1、输出长度保持不变的滑动平均操作，用于趋势提取，不用于降低时间分辨率。
随后，模块对每个样本在时间维和变量维上求全局均值，得到样本级标量特征，并由线性层和 Softmax 生成对应的门控权重：
门控权重使用样本级全局均值生成，而不是引入更复杂的注意力网络或逐时间点门控，是为了控制模块复杂度。全局均值虽然只提供较粗粒度的信息，但它能够反映样本整体水平和不同样本之间的尺度差异；线性层和 Softmax 则将这一标量特征映射为三个窗口的归一化权重。Softmax 的作用在于使不同窗口贡献具有可比较的比例关系，避免融合权重无限放大，并使最终趋势项可以被理解为多个趋势候选的加权平均。
\begin{equation}
    \boldsymbol{\alpha}^{(b)} = \mathrm{softmax}(\mathbf{W}\cdot \mathrm{Mean}_{T,C}(\mathbf{X}^{(b)})+\mathbf{b}),
\end{equation}
其中 $\boldsymbol{\alpha}^{(b)}$ 表示第 $b$ 个样本在不同窗口上的融合权重。最终趋势项为：
\begin{equation}
    \mathbf{T}^{(b)} = \sum_{j=1}^{3} \alpha_j^{(b)} \mathbf{T}_j^{(b)},
\end{equation}
季节项则由原始输入减去趋势项得到：
\begin{equation}
    \mathbf{S}^{(b)} = \mathbf{X}^{(b)} - \mathbf{T}^{(b)}.
\end{equation}

该模块的特点是结构较轻，只在原有移动平均分解基础上增加了多窗口并行和一个小型门控网络。门控层初始化时使不同窗口权重接近均匀分布，从而减少新增模块对训练初期模型状态的扰动。与固定单窗口相比，多核门控分解能够根据样本级输入特征调整不同尺度趋势估计的贡献，因此更适合处理周期结构复杂的数据。该模块的具体效果将在第 4 章通过 M2 变体进行验证。
从结构解释性看，多核门控分解仍然保留了移动平均分解的直观含义。每一个窗口都对应一种可理解的趋势估计方式，门控权重则反映当前样本对不同窗口的依赖程度。与完全由深层网络学习隐式分解相比，这种形式更容易和趋势、季节项的概念对应起来。不过，该模块也有局限。样本级全局均值只能提供整体层面的自适应，不能针对每个时间点或每个变量分别选择窗口；当序列内部不同时段的变化尺度差异很大时，单个样本级权重可能仍然不够细致。因此，本文将其定位为轻量增强，而不是对分解机制的彻底重构。

\begin{figure}[htbp]
    \centering
    \includegraphics[width=0.98\textwidth]{thesis/fig_multi_kernel_gated_decomposition.png}
    \caption{多核门控分解模块结构图}
    \label{fig:multi_gated_decomposition}
\end{figure}

\section{通道混合模块}
\esection{Channel Mixing Module}

原始 TimeMixer 默认采用通道独立策略，即将多变量时间序列中的每个变量拆分为独立序列进行处理。这种策略可以降低模型复杂度，尤其适合变量数较多的数据集。但它也带来一个问题：模型在主干结构中难以显式利用变量之间的相关性。对于气象、交通、电力等多变量数据，变量间关系往往具有实际意义，因此完全忽略这部分信息可能限制模型性能。
通道独立策略的优势在于计算路径清晰，变量数增加时不会立即导致全连接式的复杂度增长，因此适合 ECL、Traffic 等高维多变量数据集。但从数据含义看，许多变量并不是完全孤立的。气象数据中的温度、湿度、气压等指标存在物理联系；交通传感器之间可能受到道路拓扑和相邻路段状态影响；电力负荷数据也可能具有用户群体或时间规律上的关联。如果模型始终把每个变量当作独立序列处理，就可能错过这些有用信息。

本文在多尺度输入处理阶段加入通道混合模块。具体而言，模型先构建多尺度序列，再对每个尺度的序列进行归一化；若启用通道混合，则在通道独立拆分和嵌入之前，对变量维度使用一个低秩 MLP 进行残差映射：
\begin{equation}
    \tilde{\mathbf{X}} = \mathbf{X} + f_{cm}(\mathbf{X}),
\end{equation}
实现时，该 MLP 作用在变量维度 $C$ 上，即对每个时间步的变量向量进行低秩映射。
其中
\begin{equation}
    f_{cm}(\mathbf{X}) = \mathbf{W}_2 \cdot \mathrm{GELU}(\mathbf{W}_1 \mathbf{X}).
\end{equation}
具体而言，第一层线性映射将变量维度从 $C$ 降到 $\min(64,C)$，第二层再映射回原变量维度。这样可以在控制参数量的同时完成一次变量间信息交换。
通道混合层采用残差形式，目的是避免新增模块对原始输入造成过强扰动。初始化时，残差分支输出接近于零，使训练开始时该模块近似恒等映射。需要说明的是，本文通道混合只在多尺度输入进入主干前进行一次浅层交互，因此它的建模能力有限。该模块的具体效果将在第 4 章通过 M3 及相关组合变体进行分析。
\begin{figure}[htbp]
    \centering
    \includegraphics[width=0.95\textwidth]{thesis/fig_channel_mixing_module.png}
    \caption{通道混合模块结构图}
    \label{fig:channel_mixing_module}
\end{figure}

\section{模型变体与消融设计}
\esection{Model Variants and Ablation Design}

为了分析三个模块的独立贡献和组合效果，本文构造了 8 个模型变体，如表 \ref{tab:variants} 所示。其中 M0 表示 TimeMixer 基线模型，M1、M2、M3 分别表示只加入一个模块，M12、M13、M23 表示两两组合，M4 表示三个模块全部加入。

\begin{table}[htbp]
    \centering
    \caption{模型变体设置表}
    \label{tab:variants}
    \begin{tabular}{ccccc}
        \toprule
        变体 & 深度可分离下采样A & 多核门控分解B & 通道混合C & 含义 \\
        \midrule
        M0 & 否 & 否 & 否 & TimeMixer 基线 \\
        M1 & 是 & 否 & 否 & 仅替换下采样 \\
        M2 & 否 & 是 & 否 & 仅替换分解模块 \\
        M3 & 否 & 否 & 是 & 仅加入通道混合 \\
        M12 & 是 & 是 & 否 & A+B \\
        M13 & 是 & 否 & 是 & A+C \\
        M23 & 否 & 是 & 是 & B+C \\
        M4 & 是 & 是 & 是 & 三模块完整模型 \\
        \bottomrule
    \end{tabular}
\end{table}

这种设计覆盖了三个二值模块的全部组合。通过 M0 与 M1、M2、M3 的比较，可以观察单个模块独立加入后的效果；通过 M4 与 M12、M13、M23 的比较，可以观察完整模型去掉某个模块后的变化；通过所有变体的总体胜出次数，可以判断哪一类改进更稳定。本文后续实验结论显示，模块组合并不一定带来线性叠加收益，这也是本文结果分析中的一个重要发现。
